\section{Introduction}
\label{sec:intro}

Large language model (LLM) agents have crossed from answering questions to acting on the Web: they browse pages, edit documents, move money, and post content by orchestrating sequences of tool calls against external services~\cite{debenedetti2024agentdojo,bai2026harnessrisk}. This deployment shift changes what safety means. In a chatbot, a failure is a harmful utterance; in an agent, a failure is a \emph{deed}---a privileged tool invoked at the wrong time, with the wrong arguments, for the wrong reason---and deeds are composed of many individually innocuous steps. Recent multi-turn audits show that this gap between what agents can do and what they should do widens precisely as interactions lengthen and tool access deepens~\cite{li2026unsafer}.

The threat surface of the agentic Web has likewise matured into families rather than incidents. Indirect prompt injection turns any tool result, webpage, or document the agent reads into an attack channel, with automatically generated and covert payloads demonstrated at scale~\cite{wang2026landscape,wang2025webinject,liu2024automatic,lee2026will}; adversarial prompts can even steer which tools the agent selects~\cite{shi2025prompt}. Separately, malicious users can request harmful goals outright, and benchmark suites now measure agent harmfulness and safety awareness across tool stages~\cite{andriushchenko2024agentharm,ye2024toolsword,yuan2024rjudge}. A third family receives less attention but is arguably more damaging in deployment: over-privileged tool use, where an agent performs an action that the task did not warrant, without any injected payload at all~\cite{yang2026when}. Defenses have proliferated---structured queries~\cite{chen2024struq}, preference-optimized resistance~\cite{chen2024secalign}, tool-dependency-graph plan conformance~\cite{an2025ipiguard}, and step-level guardrails and diagnostic frameworks~\cite{mou2026toolsafe,liu2026agentdog,hossain2026nexus,chen2025agentguard}---yet almost all of them judge the agent one response or one step at a time.

We argue that this per-step granularity is the wrong audit unit, and that the systems-security community settled the analogous question a decade ago. Host intrusion detection moved from signature matching on individual syscalls to \emph{provenance graphs}, because attacks such as data exfiltration are invisible in any single event and emerge only from the connectivity of many events~\cite{zipperle2022provenancebased}. The same structural insight transfers to agents: an injection hijack is a dataflow edge from an untrusted observation into a sensitive sink; over-privilege is a privileged call whose authorizing subgraph is missing; a malicious agent is a workflow whose entire call structure diverges from task service. Text-level auditors cannot see any of this, because the pattern is not in the tokens---it is in the graph. Provenance-graph detectors, however, also teach a cautionary lesson: mimicry attacks that splice benign structure into attack structure defeat even graph-based detectors~\cite{goyal2023sometimes}, so any credible audit protocol must test robustness, not just accuracy~\cite{goyal2024orchid}.

Translating this insight into a benchmark surfaces three obstacles that, to our knowledge, no prior work has addressed together. First, agent traces are \emph{sequences with semantics}, and the closest existing resources model them as such: AgentDrift publishes step-labeled injection-hijacked trajectories and DriftNet detects and localizes hijacks with a trajectory transformer~\cite{pinjari2026agentdrift,pinjari2026driftnet}, while benchmark suites such as AgentSecBench, HarnessRisk, PASTABench, AutoDojo, and WAInjectBench evaluate agents or detectors on textual trajectories~\cite{alpay2026agentsecbench,bai2026harnessrisk,sun2026pastabench,ma2026autodojo,liu2025wainjectbench}. None of them constructs the provenance graph that makes structural auditing possible. Second, graph anomaly detection (GAD) has matured from node-level benchmarks~\cite{liu2022bond,tang2023gadbench} to graph-level and foundation-model methods~\cite{qiao2024deep,wang2026glass}, but its standard testbeds---citation and transaction networks---share none of the temporal, typed, small-graph structure of agent traces, so headroom on those testbeds says little about the agentic Web. Third, evaluation practice itself is under audit: prompt-injection detectors have been shown to pass the very tests they were tuned on~\cite{liu2026passing}, and agent-safety benchmarks have been criticized for weak construct validity~\cite{wang2026safety}; a new benchmark must therefore build generalization and contamination controls into its spine rather than appending them as afterthoughts.

This paper introduces TraceGraph-Bench, a benchmark and evaluation protocol that compiles agent tool-call trajectories into heterogeneous provenance graphs and treats trajectory-level auditing as graph-level anomaly detection. Our data engine couples three anomaly families under one graph schema: injection-hijacked and benign trajectories replayed from AgentDojo and AgentDrift~\cite{debenedetti2024agentdojo,pinjari2026agentdrift}, malicious-use trajectories from AgentHarm~\cite{andriushchenko2024agentharm}, and a newly constructed over-privileged tool-use family that operationalizes the taxonomy of~\cite{yang2026when}. A disclosure ablation contrasts structure-only graphs with structure-plus-text graphs, isolating what structure alone contributes. The protocol demands leave-one-attack-family-out (LOAO) generalization, cross-environment transfer between replayed and newly generated environments, robustness to mimicry-style splicing~\cite{goyal2023sometimes}, and an audit-economics accounting of detector cost per trace. Our reference detector---a supervised heterogeneous-graph detector over the compiled graphs---together with a matrix of rule-based, sequence-based, node-level GAD, and graph-level GAD baselines yields a first systematized map of this new problem, showing [TBD---headline findings preview]. Our contributions are:
\begin{itemize}
  \item \textbf{Trajectory provenance graphs.} A typed construction protocol that compiles raw agent traces into heterogeneous provenance graphs (task, model step, tool call, argument entity, result, sensitive sink; invocation, argument-flow, and temporal edges), with a disclosure knob that releases graphs at structure-only or structure-plus-text granularity.
  \item \textbf{The first graph-level auditing benchmark.} Three anomaly families under one schema, spanning injection hijack, malicious use, and over-privilege, with LOAO, cross-environment transfer, mimicry robustness, and audit-economics metrics designed against the contamination failure modes documented in prior audits.
  \item \textbf{A systematized detector matrix.} Baselines from hand-crafted rules and sequence models to repurposed node-level GAD and modern graph-level GAD, plus a supervised heterogeneous-graph reference detector, all evaluated on identical graph inputs.
  \item \textbf{Empirical findings on structural auditability.} Evidence on when structure alone suffices, when text is necessary, how detectors generalize to unseen attack families, and where mimicry breaks them [TBD---fill with concrete numbers].
\end{itemize}
Section~\ref{sec:related} situates our contribution against agent safety, GAD, and provenance research; Sections~\ref{sec:method} and~\ref{sec:bench} formalize the graph construction and benchmark protocol; Sections~\ref{sec:setup} and~\ref{sec:results} report the study; Section~\ref{sec:conclusion} concludes.
